\documentclass[12pt]{article}
\usepackage[utf8]{inputenc}
\usepackage[T1]{fontenc}
\usepackage[margin=1in]{geometry}
\usepackage{setspace}
\usepackage{xcolor}
\usepackage[colorlinks=true, urlcolor=blue!60!black, linkcolor=blue!60!black]{hyperref}
\setlength{\parindent}{2em}
\setlength{\parskip}{1.25em}
\renewcommand{\baselinestretch}{1.0}
\begin{document}
\begin{center}
{\LARGE \textit{Brown v. Board of Education (Brown II)}, 349 U.S. 294 (1955)\par}
\vspace{0.8em}
{\large Case Brief --- Four-Page Edition\par}
\end{center}
\vspace{1em}

\section{Memory Jogger}

The Court delegated primary responsibility for implementing school desegregation to local school boards, supervised by the lower federal courts using equitable flexibility, and ordered Black students admitted to public schools on a racially nondiscriminatory basis ``with all deliberate speed'' (\href{https://supreme.justia.com/cases/federal/us/349/294/}{Justia}).

\section{Detailed Case Facts}

On May 17, 1954, \textit{Brown I} and \textit{Bolling v. Sharpe} held segregated public education unconstitutional, but the Court set the cases for further argument on relief: whether to require immediate admission or permit gradual adjustment, and who should frame the decrees. The consolidated cases came from Topeka, Kansas; Clarendon County, South Carolina; Prince Edward County, Virginia; the District of Columbia; and New Castle County, Delaware. The Court invited the Attorney General and the Attorneys General of the segregating states to participate, and six states filed briefs for four days of reargument in April 1955. By then the District of Columbia, Kansas, and Delaware had taken substantial steps toward desegregation, while South Carolina and Virginia had done nothing. Southern leaders were mobilizing massive resistance, and NAACP counsel led by Thurgood Marshall urged firm, prompt deadlines.

\section{Procedural History}

Three-judge district courts denied relief in the Kansas, South Carolina, and Virginia cases, and the District of Columbia court dismissed \textit{Bolling}. Only in Delaware did the state courts order immediate admission of the Black plaintiffs, because the Black schools were materially inferior. After \textit{Brown I}, the Court retained jurisdiction for reargument on remedies. On May 31, 1955, it reversed and remanded the four federal cases and affirmed the Delaware judgment, remanding for any further proceedings.

\section{Judicial Votes}

Unanimous, \textbf{9--0} (decided May 31, 1955).
\begin{itemize}
\item \textbf{Majority:} Chief Justice Warren, joined by Justices Black, Reed, Frankfurter, Douglas, Burton, Clark, Minton, and Harlan.
\item \textbf{Concurrences and dissents:} None.
\end{itemize}

\section{Holding}

All federal, state, and local laws requiring or permitting racial discrimination in public education must yield to \textit{Brown I}. Local school authorities hold primary responsibility for desegregation plans; the original trial courts, guided by equitable principles, judge whether they act in good faith. Courts must require a ``prompt and reasonable start'' toward full compliance, place on defendants the burden of justifying additional time, retain jurisdiction, and order admission on a racially nondiscriminatory basis ``with all deliberate speed.'' The listed administrative factors are non-exhaustive guidance, not a mandatory checklist.

\section{Analysis of Opinions}

\textbf{Majority (Warren).} Having reaffirmed the constitutional right, Warren treated the remaining issue as one of institutional competence and remedy. \emph{First}, he decentralized remedial design: school authorities must solve varied local problems, and the trial courts that heard the cases, given their proximity to local conditions, must determine whether officials' actions constitute ``good faith implementation of the governing constitutional principles.'' \emph{Second}, he grounded the remedy in traditional equity, which balances the plaintiffs' personal interest in admission as soon as practicable against the public interest in removing administrative obstacles during the transition---while warning that constitutional principles may not yield to disagreement. \emph{Third}, he set a two-step burden framework: a prompt and reasonable start first, and only then additional time, which defendants must prove necessary in the public interest and consistent with good-faith compliance at the earliest practicable date. Courts could weigh the school plant, transportation, personnel, revision of attendance areas into compact units, and changes to local laws. No Justice wrote separately; Warren sought to preserve the Court's unanimity in the face of anticipated noncompliance.

\section{Examples: Future Applications}

\begin{enumerate}
\item \textbf{Same-Side: A Valid One-Year Transition.} A rural board promptly adopts a binding plan to redraw bus routes, consolidate buildings, reassign teachers without regard to race, and admit Black students next year. It has made a prompt and reasonable start and tied a brief delay to real administrative obstacles.
\item \textbf{Same-Side: Prison Reform Decree.} A court finding unconstitutional prison medical care orders a staged compliance plan within 90 days and retains jurisdiction to monitor good-faith implementation, mirroring \textit{Brown II}'s equitable model.
\item \textbf{Opposite-Side: Delay for Public Hostility.} A board seeks a three-year postponement citing public opposition and threatened unrest. Hostility is not a permissible administrative obstacle.
\item \textbf{Opposite-Side: Indefinite Tokenism.} Five years on, a district transfers two Black students through a pupil-placement test and claims it needs time to study community readiness. Without a prompt start, it cannot carry its burden.
\item \textbf{Fence-Sitter: Grade-a-Year Plan.} A large district starts with first grade now and adds a grade a year over twelve years. It begins immediately, but students already in upper grades would never attend a desegregated school.
\end{enumerate}

\section{Critique}

\textbf{Rights without remedies.} \textit{Brown II} placed primary responsibility on the very boards that had defended segregation, and its oxymoronic formula set no deadline, inviting delay. More than a decade later most Southern schools remained segregated, and \textit{Alexander v. Holmes County Board of Education} (1969) repudiated the standard. Progressive scholars such as Ogletree, Bell, and Garrow call it a tragic compromise that traded immediate justice for unanimity.

\textbf{Competing values.} The opinion pits the principle that constitutional rights are personal and present against institutional pragmatism, which fears that aggressive mandates without political backing will provoke defiance and erode the judiciary's legitimacy. By prioritizing unanimity and gradualism, the Court sacrificed a generation of Black schoolchildren's present rights to the hope that moderation would elicit voluntary compliance.

\textbf{Originalist critiques.} A plaintiff who proves a statute unconstitutional is traditionally entitled to an immediate personal remedy, not an indefinite phase-in. Justice Thomas and others fault the decision for launching open-ended structural injunctions beyond the historic powers of equity, so that it gave plaintiffs too little relief and district judges too much managerial power.

\section{Key Quotations}

\begin{quote}
``Traditionally, equity has been characterized by a practical flexibility in shaping its remedies and by a facility for adjusting and reconciling public and private needs.''
\end{quote}
\noindent --- Chief Justice Warren, for the Court, 349 U.S. at 300 (footnotes omitted).

\begin{quote}
``School authorities have the primary responsibility for elucidating, assessing, and solving these problems; courts will have to consider whether the action of school authorities constitutes good faith implementation of the governing constitutional principles.''
\end{quote}
\noindent --- Chief Justice Warren, for the Court, 349 U.S. at 299.

\begin{quote}
``But it should go without saying that the vitality of these constitutional principles cannot be allowed to yield simply because of disagreement with them.''
\end{quote}
\noindent --- Chief Justice Warren, for the Court, 349 U.S. at 300.

\end{document}
